\section*{Answer}
\begin{enumerate}\setlength\itemsep{4pt}
\item \textbf{H\'ector's YOLO request (\iss{58}): which microkernel sizes run fastest?}
  \begin{itemize}\setlength\itemsep{1pt}
  \item \textbf{1 thread:} 24x8 for the whole network: 697.5 ms per pass over the 50 shapes (4.26 GOPs, \iss{109} re-search, \S\ref{sec:best}). In the \iss{61} sweep, 24x8 698.9 ms against 785.3 ms (3.79 GOPs) for our GEMM at H\'ector's tile (16x8; 16x4 at n = 4). Per n: 24x4 (n = 4), 24x9 (n = 8--32), 24x8 (n = 64, 256), 32x7 (n = 80, 128); at n = 4, 8 and 256 the winner changes between launches (n = 256: 24x8 and 32x7 tie).
  \item \textbf{8 threads, in our parallel GEMM:} the winner changes with n and is best read as a family: MR 32 / NR 4 for n $\le$ 8, MR 32 / NR 6--7 for n = 16--64, MR 16 / NR 12--15 for n $\ge$ 80. Best single tile: 32x7, 140.6 ms per pass (\iss{109} re-search, \S\ref{sec:best}). In the \iss{63} sweep (one run): best tile per shape 134.2 ms, 32x7 145.9 ms, H\'ector's tile in our GEMM 149.8 ms.
  \item \textbf{Our GEMM against H\'ector's, one series (\iss{109}), ours at its best single tile (24x8 at 1T, 32x7 at 8T) against his unchanged code, the 50 layers with n $>$ 1 (n = 1, GEMV, left out: no GEMV in our design), same settings and thread binding:} 1T: H\'ector's 5.07 GOPs, ours 4.23 (vector packer), 4.05 (scalar packer); H\'ector's 1.20$\times$ faster. 8T: ours 21.14 against H\'ector's 19.00 GOPs, ours 1.11$\times$ faster, outside the launch range; the lead comes from Y00 (his about 29 ms, ours about 7 ms): without Y00 H\'ector's leads, 22.78 against 21.81 GOPs (\S\ref{sec:109}).
  \item \textbf{Why ours trails H\'ector's at 1 thread (\iss{95}, \S\ref{sec:95}):} our A packer reads the column-major A one 24-row panel at a time, jumping a full column per read; H\'ector's reads each column once. The pack is 77\% of our time on Y10 and 90\% on Y05 (instrumented build, shares only). The 1T gap sits on the 23 shapes with m $>$ 300 and k $\le$ kc (131.0 of 140.5 ms), where a board pack test puts the read-order cost at the same size (162.2 ms). The fix is pending in \iss{96}.
  \item \textbf{Caveat:} the tiles were timed in \emph{our} GEMM, not in H\'ector's. Without Y00, H\'ector's engine is 1.04$\times$ ahead of ours at 32x7 (\iss{109}, both bound); tile for tile (ours at 16x8, his threads not bound, \iss{63}) it was about 15\% ahead (\S\ref{sec:63}). Whether our tile family speeds up \emph{H\'ector's} GEMM is the one essential test still open (\S\ref{sec:open}).
  \end{itemize}
\item \textbf{Paper shapes (\iss{71}): does the paper's tile choice hold?} At 1 thread, yes among our 19 tiles: 32x7 wins 9 of 10 shapes (16x15 wins B2). At 8 threads, no: 16x8 wins 8 of 10 shapes and the whole pass (48.27 GOPs, against 44.07 for 32x7).
\item \textbf{Packing (\iss{83}, \iss{85}): keep the vector packer.}
  \begin{itemize}\setlength\itemsep{1pt}
  \item YOLO with a scalar packer: scalar $\div$ vector GOPs median 0.929 (1T) and 0.945 (8T), down to 0.58 at n = 4; the best tile changes on 13 (1T) and 20 (8T) of 50 layers.
  \item BERT and Square: mostly within the 3\% board/day gap. Exceptions: B3 at 1T (1.227), the shape where the vector packer is known to be slow (\iss{36}); B2 at 1T (0.961, steady); B5 at 8T (0.963) and B3 at 8T (0.941), each with one slow scalar launch.
  \item Slowdown $\approx$ pack share $\times$ extra pack cost. n sets the pack share (87\% at n = 4, 7\% at n = 1024 at k = 256); k and m set the extra cost (1.2--3.4$\times$ when the packed blocks stay in cache, 1.02--1.17$\times$ when both packers wait on memory, at m = 1200).
  \end{itemize}
\item \textbf{Across YOLO and the paper shapes, the best tile at 1 thread is not the best at 8 threads.}
\end{enumerate}

\noindent\textbf{Metric note.} Packing numbers are recomputed with the true GOPs by mean time; the \iss{83} / \iss{85} reports used a different average. The answers hold either way (\S\ref{sec:metric}).

\section*{Glossary}
{\small
\begin{tabularx}{\textwidth}{@{}>{\bfseries}p{3.6cm} X@{}}
\toprule
m $\times$ n $\times$ k & C (m $\times$ n) += A (m $\times$ k) $\cdot$ B (k $\times$ n). For YOLO and BERT, nc = n, so each packed element of A is reused n times. \\
tile MR$\times$NR & The microkernel's block of C: MR rows by NR columns. The sweeps \iss{61} / \iss{63} try 19 Candidate Tiles per shape; the \iss{109} re-search the leading 4 (1T), and the leading 7 plus 24x8 (8T). \\
cache blocking mc / nc / kc & Block sizes of the GEMM loops: mc rows of A, nc columns of B, kc depth, per packed block. Our GEMM takes them from the analytical cache model for each shape and tile (clamped to m, n, k; no nc overrides); they do not change with threads or packer. Every value quoted here is the one the run's console log prints. Chart labels give the tile, then its blocking. \\
best tile & For one shape (or one n group), the tile with the highest GOPs by mean time, averaged over launches. \S\ref{sec:1t}, \S\ref{sec:8t} and Part 2 average each shape's GOPs over launches; \S\ref{sec:best}--\S\ref{sec:94} and Part 5 average the times. The two differ by under 0.4\%, except that at 1T, n = 256, 24x8 and 32x7 swap places (both $\approx$ 7.56--7.60). \\
family & A group of tiles that trade places between launches; the gap between them is inside the launch range. \\
GOPs by mean time & 2mnk $\div$ mean time per call (raw column \texttt{KerGOPs\_Harmonic}). Main metric. \\
only GOPs by mean time & Every GOPs and ms value in this report is by mean time. At 1T the median rep agrees with it within 0.019 (ratio) on 95\% of cases; at 8T a few slow reps per launch raise the mean time, so GOPs by mean time are lower on mid-size YOLO shapes (median ratio 1.02--1.12, \texttt{metric\_gap.py}). \\
ms per pass & Sum of one call's mean time over the shapes (each distinct shape once, unweighted). \\
best single tile & The one tile, at a given thread count, with the shortest pass over all layers; used for every layer (user's rule, \S\ref{sec:best}). \\
re-search & \iss{109}: the leading tiles of \iss{61} / \iss{63} timed again on the YOLO layers with more launches (ranked over the 50 with n $>$ 1), to settle the best single tile. \\
launch & One full run of the test (every tile on every shape). Written ``launch 1/2/3''. \\
A-B-B-A & Run order that interleaves two sides (e.g.\ ours, H\'ector's, H\'ector's, ours), so slow drift hits both. \\
launch range / $\sim$ & Lowest to highest launch. A gap smaller than the wider range is not called a difference (marked $\sim$, or an open dot). \\
1T / 8T & 1 thread pinned to core 4 / 8 threads on cores 0--7 (core 0 only for 8-thread runs, by the user's exception). \\
V, S, H & In \S\ref{sec:109} (\iss{109}) and \S\ref{sec:94} (\iss{94}): V our GEMM with the vector packer, S with the scalar packer, H H\'ector's \texttt{gemm.bin} (his code unchanged). \\
thread binding & \texttt{OMP\_PROC\_BIND=close OMP\_PLACES=cores}: each OpenMP thread stays on one core. Ours at 8T always; H\'ector's at 8T only in \iss{94} and \iss{109}. \\
uninstrumented build & \texttt{INSTRUMENTED=0}, no timers inside the GEMM. All GOPs come from it. \\
instrumented build & Timers around pack and compute (\texttt{--instr}). Timers cost 11--17\%, so only its time shares and ratios are used. \\
reference protocol & No warm-up; reps repeat until \texttt{--tmin} seconds (2 s; 3 s in \iss{85}; 1 s in the \iss{95} probe runs) have passed. \\
warm-up protocol & Warm-up first, then a fixed 200 reps. Used by the instrumented build. \\
vector / scalar packer & Copies a block of A or B into a contiguous buffer before the microkernel reads it. Vector: RVV loads/stores, 8 fp32 per instruction. Scalar (\texttt{--scalar-pack}): one float per load/store. Same layout and buffers. \\
pack share & Pack time $\div$ (pack + compute), from the instrumented build; A and B together. \\
extra pack cost & Scalar pack time $\div$ vector pack time. \\
s/v & Scalar GOPs $\div$ vector GOPs. Below 1: vector packing is faster. \\
predicted s/v & (vector pack + compute) $\div$ (scalar pack + compute): the slowdown if only the pack time changes. \\
A, lda & A is m $\times$ k, stored column by column; the next column starts lda = m floats later. \\
panel-outer / column-outer & Two orders to pack A into MR-row panels. Panel-outer (ours): for each panel, for each of the kc columns, read MR floats. Column-outer (H\'ector's): for each column, for each panel, read MR floats. Same bytes, same packed result. \\
probe run & One shape, ours then H\'ector's, back to back, 1T (\texttt{probe\_issue95.sh}, \texttt{--tmin 1}). \\
pack test & \texttt{issue95/pack\_order.c}: packs the same A both ways on the board and times each. Reads high in absolute terms; it gives the size of the order effect. \\
order cost / gap & Pack time in our order minus in H\'ector's order (same blocks); our time minus H\'ector's time per call at 1T (\iss{94}). \\
oracle & The naive scalar GEMM that checks C in correctness runs. \\
\bottomrule
\end{tabularx}}

\section*{Setup}
{\small
\begin{tabularx}{\textwidth}{@{}>{\bfseries}p{2.6cm} X@{}}
\toprule
chip & SpacemiT K1: 8 X60 cores, RVV 1.0, VLEN 256. Same chip as H\'ector's OrangePi. \\
boards & \texttt{jbpi} and \texttt{jbpi2}, reached through a jump host. Each table names its board. \\
house rules & Idle board (other users checked before every launch), \texttt{--no-nc-overrides}, every quoted number from a committed log. Usually 3 launches per side, interleaved (A-B-B-A or rotated order); each table states its own launches and order. \\
shapes & YOLO26: H\'ector's 53 GEMM shapes after im2col: \texttt{datasets/yolo26\_k1.dat} (50, n $>$ 1) and \texttt{yolo26\_k1\_n1.dat} (3, n = 1; together \texttt{yolo26\_k1\_53.dat}). Our GEMM times the 50 in \iss{61}, \iss{63}, \iss{94} and all 53 in \iss{109}; H\'ector's \texttt{gemm.bin} always runs all 53 (n = 1 on his GEMV). Our design has no GEMV, so every result here uses the 50 layers with n $>$ 1. Paper: Table 1 of arXiv 2603.17800v1 (\texttt{datasets/paper\_bert\_square.dat}). \\
H\'ector's GEMM & GEMMLib\_yolo at commit \texttt{b292c92}, unedited, built on the board with H\'ector's Makefile; \texttt{gemm.bin}, \texttt{TIMIN=2.0}; H\'ector's driver reports the mean time only. \\
\bottomrule
\end{tabularx}}
